\section{Homeomorphism groups of zero-dimensional spaces}\label{sec:zero-dim}

For the Cantor set, and more generally for every compact zero-dimensional
space whose nonempty clopen subsets are all homeomorphic, the flow of maximal
chains is the universal minimal flow of the homeomorphism group.

\begin{definition}[Proximal flows]\label{def:proximal}
\uses{def:flow}
Two points $x,y$ of a $G$-flow $X$ are \emph{proximal} if there is a net
$(g_i)$ in $G$ such that $(g_ix)$ and $(g_iy)$ converge to the same point. The
flow is \emph{proximal} if every pair of its points is proximal. A factor of
a proximal flow is proximal.
\end{definition}

\begin{theorem}[Glasner--Weiss \cite{GlasnerWeiss2003}]\label{thm:gw-cantor}
\uses{prop:maximal-chains, def:umf, def:proximal, thm:ellis}
Let $E$ be the Cantor set and let $G=\mathrm{Homeo}(E)$, with the topology of
uniform convergence. Then $M(G)$ is isomorphic to the flow $\Phi(E)$ of
maximal chains in $E$. In particular $M(G)$ is metrizable, and in fact
homeomorphic to the Cantor set. Moreover, $\Phi(E)$ is proximal, so every
minimal $G$-flow is proximal.
\end{theorem}

\begin{example}[The standard chain]\label{ex:standard-chain}
\uses{thm:gw-cantor, thm:pestov-ea}
Realise $E$ as the ternary Cantor set in $[0,1]$. The chain
$c_0=\{E\cap[0,t] : t\in[0,1]\}$ is maximal, with root $0$. Its stabiliser is
the group $P$ of order-preserving homeomorphisms of $E$, and $P$ is isomorphic
to $\mathrm{Aut}(\mathbb{Q},<)$, through its action on the complementary open
intervals of $E$, which are ordered like the rationals. Hence $P$ is extremely
amenable by Theorem~\ref{thm:pestov-ea}. Moreover, $P$ has exactly six fixed
points in $\Phi(E)$, and the action of $G$ on $\Phi(E)$ is minimal
\cite[Theorems~1.2 and~1.3]{GlasnerWeiss2003}.
\end{example}

\begin{proof}[Proof of proximality, after Glasner and Weiss]
Let $c_1,c_2\in\Phi(E)$ and let $Z$ be the closure of the orbit of
$(c_1,c_2)$ in $\Phi(E)\times\Phi(E)$. Since $P$ is extremely amenable, $Z$
contains a pair $(a_1,a_2)$ of $P$-fixed points, and Glasner and Weiss check
directly that any two of the six $P$-fixed chains are proximal. So $Z$
contains a diagonal point $(a,a)$, which means that $c_1$ and $c_2$ are
proximal.
\end{proof}

\begin{remark}[Stone duality]\label{rem:stone-duality}
\uses{thm:gw-cantor, ex:kpt-computations, thm:ellis}
The clopen subsets of $E$ form the countable atomless Boolean algebra
$\mathbf{B}_\infty$, and $\mathrm{Homeo}(E)$ is isomorphic, as a topological
group, to $\mathrm{Aut}(\mathbf{B}_\infty)$ with the topology of pointwise
convergence. Kechris, Pestov and Todorcevic computed the universal minimal flow
of $\mathrm{Aut}(\mathbf{B}_\infty)$ as the space of natural orders on
$\mathbf{B}_\infty$ \cite[Section~8]{KPT2005}
(Example~\ref{ex:kpt-computations}). By Theorem~\ref{thm:ellis}, this flow is
isomorphic to $\Phi(E)$.
\end{remark}

\begin{definition}[h-homogeneous spaces]\label{def:h-homogeneous}
A compact zero-dimensional space $X$ is \emph{h-homogeneous} if all its
nonempty clopen subsets are homeomorphic to each other; equivalently, by Stone
duality, if its algebra of clopen sets is a homogeneous Boolean algebra. A
metrizable h-homogeneous space with more than one point is homeomorphic to the
Cantor set.
\end{definition}

\begin{theorem}[Glasner--Gutman \cite{GlasnerGutman2012}]\label{thm:glasner-gutman}
\uses{def:h-homogeneous, prop:maximal-chains, thm:ellis}
Let $X$ be an h-homogeneous compact zero-dimensional space and let
$G=\mathrm{Homeo}(X)$, with the compact-open topology. Then $M(G)$ is isomorphic
to $\Phi(X)$.
\end{theorem}

\begin{remark}\label{rem:glasner-gutman}
\uses{thm:glasner-gutman}
For metrizable $X$ this is Theorem~\ref{thm:gw-cantor}, although the proof of
Glasner and Gutman is new even in that case, since it uses no fixed point
theorem. As for the Cantor set, the main combinatorial tool is the dual Ramsey theorem,
a consequence of the
Graham--Rothschild theorem on $n$-parameter sets \cite{GlasnerWeiss2003}.
Non-metrizable h-homogeneous spaces give many new examples, among them the
generalised Cantor cubes $\{0,1\}^\kappa$ for uncountable $\kappa$ and the
corona $\beta\omega\setminus\omega$ \cite{GlasnerGutman2012}.
\end{remark}
